\documentclass{article}
\usepackage{graphicx} % Required for inserting images
\usepackage{hyperref}
\usepackage{parskip}

\title{Momoni - Vägledning för handledare}
\date{}

\begin{document}

\maketitle

\section{Välkommen till Momoni}
Välkommen som handledare för Cyber Threat Challenge – Momoni. 

I detta dokument hittar du allt du behöver för att genomföra passet. Upplägget är interaktivt och bygger på att deltagarna kliver in i rollen som cybersäkerhetsspecialister i det fiktiva landet Momoni. 

Rollspelet börjar direkt när passet startar. 

Materialet har tagits fram av \textit{Einar Bergviken} och \textit{Nils Sjöstrand Berg}. 

Vid frågor kan ni kontakta: 
 Einar Bergviken - \hyperlink{einar.bergviken@ungaforskare.se}{einar.bergviken@ungaforskare.se} 
 Nils Sjöstrand Berg - \hyperlink{nils.sjostrand.berg@ungaforskare.se}{nils.sjostrand.berg@ungaforskare.se} 

\subsection{Snabb överblick}
\textbf{Antal handledare:} 2-3 handledare
\textbf{Längd:} 2 timmar
\textbf{Upplägg:} Interaktivt rollspel i tre akter
\textbf{Grupper:} Upp till 6 myndigheter/grupper (minst 3 och upp till 6 deltagare per grupp)
\textbf{Pauser:} Ta en 5-10 minuters paus mellan varje akt
\textbf{Material:} I filen \textit{material} som ligger i samma mapp som detta dokument finns en detaljerad beskrivning av allt material som ska skrivas ut.

\textbf{Syfte}
Deltagarna ska:
\begin{itemize}
    \item Förstå vanliga cyberhot
    \item Göra riskbedömningar
    \item Klassificera information
    \item Fatta beslut under begränsade resurser
    \item Reflektera kring ansvar och konsekvenser
\end{itemize}
Handledarens uppgift är att driva berättelsen framåt, ställa frågor, förtydliga begrepp och skapa engagemang; inte bara att ge “rätta svar”.

\tableofcontents

\section{Kort sammanfattning av passet}
Passet utspelar sig i det fiktiva landet Momoni, med huvudstad Manamon. 

Deltagarna delas in i grupper som representerar olika myndigheter. Varje grupp arbetar som cybersäkerhetsspecialister på sin myndighet. 
\subsection{Uppdelning}
Passet är uppdelat i fyra delar: 
\begin{itemize}
    \item Inledning – Uppdraget börjar 
    \item Akt 1 – Hotbild och prioriteringar 
    \item Akt 2 – Cyberkrisen 
    \item Akt 3 – Konsekvenser och ansvar 
\end{itemize}
\subsection{Olika typer av moment}
Under passet arbetar deltagarna med två typer av moment: 
\begin{itemize}
    \item Rollspel med vägval (beslutsfattande och diskussion) 
    \item Praktiska övningar (riskanalys, klassificering och incidenthantering) 
\end{itemize}
Allt material som används finns i Cybersecurity Academys Google Drive och består av både fysiska dokument och digitala filer.
\section{Inledning (25 minuter)}
Rollspelet startar här.

Deltagarna är nyanställda cybersäkerhetsspecialister som har kallats till Momoni för att stärka landets digitala säkerhet.

\textbf{Handledarens uppgift:}
\begin{itemize}
    \item Introducera situationen via PowerPoint (i roll).
    \item Skapa en känsla av allvar och uppdrag.
    \item Dela in deltagarna i grupper (3–6 per grupp).
    \item Dela ut material.
\end{itemize}
\subsection{Moment under inledning}
\textbf{1. Skriva på anställningskontrakt}
\textit{Material: Anställningskontrakt}

Grupperna skriver under sitt kontrakt.
Kontraktet fungerar senare som stöd i deras beslutsfattande, exempelvis i situationer där de behöver motivera varför de väljer att följa eller inte följa en chefs instruktion.

Det markerar att de nu officiellt är i tjänst.

\textbf{2. Läsa myndighetsinformation}
\textit{Material: Myndighetsinformation}

Grupperna läser om sin tilldelade myndighet:
\begin{itemize}
    \item Myndighetens uppdrag
    \item Kort bakgrund
    \item Den centrala applikation som är avgörande för verksamheten
\end{itemize}
Den centrala applikationen blir viktig senare i passet.

\textbf{3. Lokalisera sitt kontor}
\textit{Material: Storyhäfte (karta finns även i PowerPoint)}

Grupperna hittar sin myndighet på kartan över Manamon.
Detta förstärker känslan av att de är placerade i en verklig miljö.
\section{Akt 1 (30 minuter)}
Nu börjar deltagarnas första uppdrag på myndigheten.

Deras uppgift är att analysera hotbilden mot sin egen myndighet och ta ställning till hur verksamhetens information bör skyddas.

Akt 1 innehåller både en praktisk riskanalys och ett första vägval i rollspelet.

\subsection{Översikt av Akt 1}

Akt 1 består av tre moment:
\begin{itemize}
    \item Identifiering av vanliga attackvektorer
    \item Klassificering av information
    \item Mötet med generaldirektören (vägval)
\end{itemize}

\subsection{1. Identifiera attackvektorer och hotaktörer}

\textit{Material: Myndighetsinformation, Säkerhetspolisens riskanalys (digital), Attackinformationsblad, Digitalt arbetshäfte}

Grupperna ska gå igenom Säkerhetspolisens rapport och informationen om den tilldelade myndigheten. Utifrån detta ska de identifiera två relevanta attackvektorer som kan riktas mot myndigheten.

Med hjälp av den information som finns i Säkerhetspolisens data ska de dessutom identifiera vilken hotaktör som utför attacken samt attackens sannolikhet.

Slutligen ska grupperna beskriva hur en sådan attack skulle kunna påverka myndighetens centrala applikation samt den information som hanteras där. 

Grupperna dokumenterar sina slutsatser i det digitala arbetshäftet.

\subsection{2. Klassificering av information}

\textit{Material: Myndighetsinformation, Digitalt arbetshäfte}

Grupperna ska nu arbeta vidare med informationen om myndighetens centrala applikation. Uppgiften är att identifiera vilken typ av information som applikationen hanterar och beskriva denna kortfattat i det digitala arbetshäftet.

Utifrån detta ska grupperna klassificera informationen enligt tre informationssäkerhetsperspektiv: \textbf{konfidentialitet}, \textbf{riktighet} och \textbf{tillgänglighet}. För varje perspektiv ska de ange vilken nivå som är lämplig samt kort motivera sitt val i arbetshäftet.

I sina motiveringar ska grupperna använda Säkerhetspolisens beskrivning av hotbilden för att resonera kring vilka konsekvenser det kan få om informationen exempelvis skulle läcka, förändras eller inte vara tillgänglig.

Syftet med momentet är att grupperna ska koppla samman de hot och attackvektorer de tidigare identifierat med vilket skydd informationen behöver.

\subsection{3. Mötet med generaldirektören}

När en grupp är färdig med sin analys kallas de till ett möte. En handledare spelar rollen som generaldirektör över myndigheten, och om många grupper blir klara samtidigt kan ytterligare handledare delta som biträdande chef (eller liknande) för att hålla flera möten parallellt. Mötet hålls gärna i ett separat rum för att förstärka rollspelet.

Under mötet ska gruppen presentera sitt arbete. De ska redogöra för de attackvektorer och hotaktörer de identifierat, beskriva sina klassificeringar och förklara sina rekommendationer. Handledaren i rollen som generaldirektör bör ställa kritiska frågor, be gruppen motivera sina prioriteringar och visa särskilt intresse för kostnadsaspekten.

Efter presentationen ska generaldirektören uttrycka oro över budgeten och ifrågasätta gruppens klassificeringar. Exempelvis kan generaldirektören säga:  
\textit{“Det här ser kostsamt ut. Är det verkligen nödvändigt med så hög klassificering? Kan vi inte lägga oss på en lägre nivå för att spara pengar?”}  

Gruppen får därefter återvända till sitt arbetsrum för att diskutera situationen och bestämma hur de vill gå vidare. Handledaren uppmuntrar gruppen att reflektera över om de vill stå fast vid sina ursprungliga klassificeringar eller justera dem.

\subsection{Vägval 1 – Säkerhet eller kostnad}

Grupperna ska nu ta ställning till om de vill:
\begin{itemize}
    \item Sänka sina klassificeringar för att minska kostnaderna, eller
    \item Stå fast vid sin ursprungliga bedömning.
\end{itemize}

De ska dokumentera sitt beslut och skriva en kort motivering i sitt arbetshäfte.

Detta beslut sparas och kommer att få betydelse i Akt 3.

\section{Akt 2 (50 minuter)}

En tid efter att grupperna har fattat sina beslut i Akt 1 inträffar en cyberkris i Momoni.

Flera myndigheter utsätts för attacker och situationen utvecklas snabbt.

Akt 2 fokuserar på praktisk incidenthantering, analys och beslutsfattande under press.

\subsection{Översikt av Akt 2}

Akt 2 består av tre moment:
\begin{itemize}
    \item Presentation av incident
    \item Incidenthantering
    \item Muntlig redovisning och förebyggande åtgärder
\end{itemize}

\subsection{1. Presentation av incident}

\textit{Material: Incidentrapporter, Digitala övningsfiler}

En kort tid efter grupperna har fattat sina beslut avbryter handledaren och meddelar att en incident har inträffat. Exempelmanus för detta finns i powerpointen.

Grupperna börjar med de två incidentrapporter som är kopplade till de attacker som påverkar just deras myndighet. När dessa är färdiga fortsätter grupperna med resterande attacker och rapporter.

\subsection{2. Incidenthantering}

\textit{Material: Digitala övningsfiler}

Grupperna ska arbeta med de digitala övningsfiler som finns i alla mappar under "Övningsmaterial"-mappen i Akt 2. För varje incident ska de analysera attacken, bedöma hur den påverkar myndigheten och föreslå lämpliga åtgärder.

\subsection{3. Muntlig redovisning}

När en grupp har hanterat en incident ska de redovisa för en handledare.

Redovisningen ska innehålla:
\begin{itemize}
    \item Vad som har hänt.
    \item Vilken attackmetod som användes.
    \item Hur myndigheten påverkades.
    \item Vilka åtgärder gruppen vidtog.
\end{itemize}

Handledaren kan ställa följdfrågor för att testa gruppens förståelse och resonemang.

\section{Akt 3 (20 minuter)}

Krisen är hanterad. Nu är det dags att stanna upp och förstå vad som egentligen hände och vad konsekvenserna blev av de beslut som fattades längs vägen.

Akt 3 kopplar ihop rollspelet med verkligheten och ger deltagarna möjlighet att reflektera kring ansvar, etik och konsekvenser av sina egna vägval.

\subsection{Översikt av Akt 3}

Akt 3 består av tre moment:
\begin{itemize}
    \item Vägval 2 – Politikern och journalisterna
    \item Diskussion i helklass
    \item Avslutning och koppling till verkligheten
\end{itemize}

\subsection{1. Vägval 2 – Politikern och journalisterna}

\textit{Material: Storyhäfte, Arbetshäfte (vägval 1 från Akt 1)}

Grupperna återvänder till sina kontor efter krisen. I storyhäftet får de läsa att journalister har samlats utanför byggnaden och att politikern (politikern är en central karaktär i storyn och finns med i storyhäftena) knackar på dörren med instruktioner om hur de ska hantera situationen.

Beroende på vilket beslut gruppen fattade i vägval 1 möter de två olika scenarier:

\subsubsection{Vägval 2.1 – Grupper som stod fast vid sin klassificering}

Politikern uppmanar gruppen att skylla cyberattacken på presidenten inför journalisterna.

Gruppen väljer ett av tre alternativ:
\begin{itemize}
    \item \textbf{Alternativ A.} Skylla ifrån sig på presidenten när journalisterna frågar.
    \item \textbf{Alternativ B.} Berätta för journalisterna hur attacken kunde hända, utan att peka ut någon specifik person.
    \item \textbf{Alternativ C.} Säga ingenting alls till journalisterna.
\end{itemize}

\subsubsection{Vägval 2.2 – Grupper som sänkte sin klassificering}

Politikern ber gruppen att inte avslöja att klassificeringen sänktes av budgetskäl, för att skydda myndighetens rykte.

Gruppen väljer ett av tre alternativ:
\begin{itemize}
    \item \textbf{Alternativ A.} Göra som politikern säger och skydda myndighetens rykte genom att ljuga för journalisterna.
    \item \textbf{Alternativ B.} Berätta för journalisterna att politikern sa åt dem att klassificera tillgänglighet lågt för att spara pengar.
    \item \textbf{Alternativ C.} Säga ingenting alls till journalisterna.
\end{itemize}

Grupperna diskuterar inom sig vilket alternativ de väljer och motiverar sitt beslut. Inget svar behöver lämnas in - diskussionen är underlaget för helklasssamtalet som följer.

\subsection{2. Diskussion i helklass}

\textit{Material: PowerPoint}

Handledarna samlar hela gruppen och leder en gemensam diskussion med utgångspunkt i rollspelet.

Förslag på diskussionsfrågor:
\begin{itemize}
    \item Vilket alternativ valde ni och varför kändes det rätt?
    \item Är det alltid rätt att följa sin chefs eller sin politikers instruktioner? Var går gränsen?
    \item Vem bär ansvaret när ett säkerhetsbeslut fattas av ekonomiska skäl och något sedan går fel?
    \item Vad skulle du ha gjort om du verkligen jobbade på en myndighet i den här situationen?
\end{itemize}

Handledaren uppmuntrar alla grupper att dela med sig av sina perspektiv, även om de fattade olika beslut. Poängen är inte att det finns ett rätt svar, utan att synliggöra att verkliga säkerhetsbeslut sällan är enkla.

\subsection{3. Avslutning – Från Momoni till verkligheten}

\textit{Material: PowerPoint}

Handledaren avslutar passet genom att lyfta ut berättelsen ur det fiktiva Momoni och koppla den till hur cyberattacker faktiskt drabbar myndigheter och organisationer i Sverige och världen.

Lyft gärna fram konkreta exempel, exempelvis:
\begin{itemize}
    \item Ransomware-attacker mot svenska kommuner och sjukhus.
    \item Statsunderstödda hackergrupper som riktar in sig på kritisk infrastruktur.
    \item Verkliga fall där bristande klassificering eller underbemanning inom IT-säkerhet fått allvarliga konsekvenser.
\end{itemize}

Avsluta med att fråga deltagarna hur passet kändes:

\textit{"Vad tar ni med er härifrån? Vad var svårast att fatta beslut om, och varför?"}

\end{document}

